\chapter{Esperanzas condicionales y opciones bermudas}

\section{Representación \guia{Teorema 14.1}}
Si $\ip{\D F}{u}=1$ y $\ip{\D G}{u}=0$,
\[
  \E[\Phi(G)\mid F=x]=\frac{\E[\Phi(G)\ind_{\{F>x\}}\dive(u)]}{\E[\ind_{\{F>x\}}\dive(u)]}.
\]
En Black--Scholes con $F=S_s$, $G=S_t$, $0<s<t$ \guia{Ejemplo 14.2}:
\[
  \dive(u)=\frac{1}{\sigma S_s}\Big(\frac{tW_s-sW_t}{s(t-s)}+\sigma\Big).
\]

\paragraph{Localización.} Con $\psi(z)=\tfrac\eta2e^{-\eta|z|}$ y primitiva $\Psi$,
\[
  \E[\Phi(G)\delta_x(F)]=\E[\Phi(G)\psi(F-x)]+\E\big[\Phi(G)(\ind_{\{F>x\}}-\Psi(F-x))\dive(u)\big],
  \qquad \eta^2=\frac{\E[\Phi(G)^2\dive(u)^2]}{\E[\Phi(G)^2]}.
\]

\section{Put bermuda por Monte Carlo \guia{§14.2}}
Programación dinámica en $0=t_0<\dots<t_n=T$:
$V_{t_n}=\Phi(S_{t_n})$, $V_{t_k}=\max\{\Phi(S_{t_k}),e^{-r\Delta t}\E[V_{t_{k+1}}\mid S_{t_k}]\}$,
evaluando la esperanza condicional en cada trayectoria con el cociente anterior
(coste $O(N^2)$ por fecha). Referencia: árbol binomial con ejercicio en las mismas fechas.

\paragraph{Fiabilidad del cociente.} El denominador estima $p_F(x)$ y es muy ruidoso en las colas
de $F$. El código estima el error estándar del cociente $R=\bar a/\bar b$ por el método delta,
$\mathrm{se}(R)=\mathrm{std}(a_i-Rb_i)/(\sqrt N\,|\bar b|)$, y devuelve \codigo{NaN} cuando el error
relativo supera un umbral (y, sin localizar, cuando hay menos de 100 muestras a cada lado de $x$),
en lugar de valores sin sentido.

\paragraph{Dos estimadores.} Tomar el máximo de estimaciones ruidosas sobre las mismas trayectorias
sesga el precio \emph{por valores} al alza (desigualdad de Jensen) y lo hace muy variable. Por eso se
calcula también el valor de la \emph{regla de ejercicio} aprendida, aplicada a trayectorias nuevas e
independientes: al ser una regla de parada admisible, da una cota inferior insesgada. Se reduce su
varianza con la put europea como variable de control.

\begin{observacion}
Con $N=2000$ trayectorias de entrenamiento y 10 fechas (12 semillas): por valores
$6{,}06\pm0{,}10$; regla de ejercicio $5{,}99\pm0{,}03$; árbol $6{,}033$. La cota inferior es
precisa y la regla aprendida es casi óptima (pérdida $<1\%$).
\end{observacion}

\begin{figure}[h]
  \centering
  \includegraphics[width=\linewidth]{08_esperanza_condicional_bermudas}
  \caption{Esperanza condicional sin regresión y precio de una put bermuda
  (\codigo{examples/08\_esperanza\_condicional\_bermudas.py}).}
\end{figure}
